% ==========================================================
\section{Research Gap}
\label{sec:research_gap}
% ==========================================================
 
The preceding survey reveals three critical gaps that emerge precisely where Behavioral Programming, feature modeling, and graphical editor frameworks intersect.
While each domain has matured significantly in isolation, the integration of formal architectural specification with visual modeling and runtime comprehensibility has received limited attention.
The following open problems, derived directly from the state-of-the-art analysis, motivate the research contributions presented in this thesis.
 
\textbf{Extensible GLSP architecture for UVL.}
BigUML demonstrates that \gls{glsp} can support multiple diagram types within a single tool, but achieves this by bundling all diagram types into one monolithic extension.
As a result, individual diagram types cannot be developed or deployed independently of one another.
How to design a modular \gls{glsp}-based architecture that permits a core diagram extension to be augmented by separately developed domain-specific extensions remains unexplored.
 
\textbf{Graphical tooling for extended UVL.}
The graphical tools currently available for \gls{uvl}, such as \texttt{flamapy.ide} and FeatureIDE, are built around the standard feature-modeling use case.
Domain-specific extensions to \gls{uvl}, such as those introduced by FMBP for Behavioral Programming, add new language constructs, types, and semantic constraints that go beyond what these general-purpose tools support.
Whether such extensions can be integrated into existing tooling without causing incompatibilities with the broader \gls{uvl} ecosystem remains an open question.
 
\textbf{Runtime integration of FMBP.}
FMBP extends \gls{uvl} to provide formal architectural comprehension of Behavioral Programming at the design phase, a valuable contribution addressed by none of the runtime comprehensibility tools surveyed in Section~\ref{sec:sota_bp_comprehension}.
However, FMBP in its current form operates entirely at the textual level and provides no visual representation of the architectural models.
Moreover, the framework does not employ the \textit{models@run.time} paradigm to synchronize the static architectural model with the live execution state of a running Behavioral Program.
Therefore, the gap between the static architectural model and the dynamic behavior of the running program remains open.

\textbf{Runtime integration of GLSP for model synchronization.}
Existing \gls{glsp}-based collaboration work focuses on enabling multiple users to edit the same diagram in real time by proxying interactive client edits through a host.
This real-time collaborative approach does not apply to the intended use case of programmatic runtime model synchronization, because it relies on interactive client-proxied edits rather than automated, server-side synchronization with a running system.
Consequently, it remains to be investigated how \gls{glsp}'s protocol and server architecture can support runtime integration patterns that enable continuous synchronization between \gls{uvl} models and live execution state.

\textbf{Thesis Contributions.}
This thesis proposes a modular, extensible \gls{glsp}-based graphical editor for the extended \gls{uvl} together with a runtime integration mechanism to explore synchronization between \gls{uvl} models and executing Behavioral Programs.
In the following chapters, the process of designing, implementing, and evaluating this tool will be addressed.